% Shared preamble for the standalone TikZ figures of the paper.
% Structure taken from the QSL slides (tizes/slides/2026-10-30-QSL/figures):
% each figure is its own document,
%   \documentclass[tikz,border=0pt]{standalone}
%   % Shared preamble for the standalone TikZ figures of the paper.
% Structure taken from the QSL slides (tizes/slides/2026-10-30-QSL/figures):
% each figure is its own document,
%   \documentclass[tikz,border=0pt]{standalone}
%   % Shared preamble for the standalone TikZ figures of the paper.
% Structure taken from the QSL slides (tizes/slides/2026-10-30-QSL/figures):
% each figure is its own document,
%   \documentclass[tikz,border=0pt]{standalone}
%   % Shared preamble for the standalone TikZ figures of the paper.
% Structure taken from the QSL slides (tizes/slides/2026-10-30-QSL/figures):
% each figure is its own document,
%   \documentclass[tikz,border=0pt]{standalone}
%   \input{figpreamble}
% then the tikzpicture. Compile it from paper/figures/ (pdflatex <name>.tex);
% the PDF lands next to it and main.tex includes it as before.
%
% Only what the figures share lives here: fonts, packages, notation.
% Each figure keeps the colours of its submitted (Inkscape) version and
% defines them itself.
\usepackage[T1]{fontenc}
\usepackage{lmodern}
\usepackage{amsmath}
\usepackage{amssymb}
\usetikzlibrary{arrows.meta, calc, fit, backgrounds}

% Fonts. The submitted figures use one of two looks, so each figure picks
% its own: serif (Latin Modern, the default here) or, by calling
% \figsansfonts in its preamble, sans text with sans Latin letters and
% digits in math (as beamer does, in the scalable Latin Modern Sans),
% Greek left in cmmi.
\def\sfletter#1{\ifx#1\relax\else
  \DeclareMathSymbol{#1}{\mathalpha}{sfletters}{`#1}\expandafter\sfletter\fi}
\newcommand{\figsansfonts}{%
  \renewcommand{\familydefault}{\sfdefault}%
  \DeclareSymbolFont{sfletters}{OT1}{lmss}{m}{sl}%
  \SetSymbolFont{operators}{normal}{OT1}{lmss}{m}{n}%
  \sfletter abcdefghijklmnopqrstuvwxyzABCDEFGHIJKLMNOPQRSTUVWXYZ\relax}

% Notation of main.tex.
\newcommand{\ket}[1]{\left|#1\right\rangle}
\newcommand{\Pauli}{\mathcal P}
\newcommand{\Clifford}{\mathcal C}

% A font size in pt, for text matched to the submitted figures.
\newcommand{\figsize}[1]{\fontsize{#1}{1.2\dimexpr#1pt\relax}\selectfont}

% then the tikzpicture. Compile it from paper/figures/ (pdflatex <name>.tex);
% the PDF lands next to it and main.tex includes it as before.
%
% Only what the figures share lives here: fonts, packages, notation.
% Each figure keeps the colours of its submitted (Inkscape) version and
% defines them itself.
\usepackage[T1]{fontenc}
\usepackage{lmodern}
\usepackage{amsmath}
\usepackage{amssymb}
\usetikzlibrary{arrows.meta, calc, fit, backgrounds}

% Fonts. The submitted figures use one of two looks, so each figure picks
% its own: serif (Latin Modern, the default here) or, by calling
% \figsansfonts in its preamble, sans text with sans Latin letters and
% digits in math (as beamer does, in the scalable Latin Modern Sans),
% Greek left in cmmi.
\def\sfletter#1{\ifx#1\relax\else
  \DeclareMathSymbol{#1}{\mathalpha}{sfletters}{`#1}\expandafter\sfletter\fi}
\newcommand{\figsansfonts}{%
  \renewcommand{\familydefault}{\sfdefault}%
  \DeclareSymbolFont{sfletters}{OT1}{lmss}{m}{sl}%
  \SetSymbolFont{operators}{normal}{OT1}{lmss}{m}{n}%
  \sfletter abcdefghijklmnopqrstuvwxyzABCDEFGHIJKLMNOPQRSTUVWXYZ\relax}

% Notation of main.tex.
\newcommand{\ket}[1]{\left|#1\right\rangle}
\newcommand{\Pauli}{\mathcal P}
\newcommand{\Clifford}{\mathcal C}

% A font size in pt, for text matched to the submitted figures.
\newcommand{\figsize}[1]{\fontsize{#1}{1.2\dimexpr#1pt\relax}\selectfont}

% then the tikzpicture. Compile it from paper/figures/ (pdflatex <name>.tex);
% the PDF lands next to it and main.tex includes it as before.
%
% Only what the figures share lives here: fonts, packages, notation.
% Each figure keeps the colours of its submitted (Inkscape) version and
% defines them itself.
\usepackage[T1]{fontenc}
\usepackage{lmodern}
\usepackage{amsmath}
\usepackage{amssymb}
\usetikzlibrary{arrows.meta, calc, fit, backgrounds}

% Fonts. The submitted figures use one of two looks, so each figure picks
% its own: serif (Latin Modern, the default here) or, by calling
% \figsansfonts in its preamble, sans text with sans Latin letters and
% digits in math (as beamer does, in the scalable Latin Modern Sans),
% Greek left in cmmi.
\def\sfletter#1{\ifx#1\relax\else
  \DeclareMathSymbol{#1}{\mathalpha}{sfletters}{`#1}\expandafter\sfletter\fi}
\newcommand{\figsansfonts}{%
  \renewcommand{\familydefault}{\sfdefault}%
  \DeclareSymbolFont{sfletters}{OT1}{lmss}{m}{sl}%
  \SetSymbolFont{operators}{normal}{OT1}{lmss}{m}{n}%
  \sfletter abcdefghijklmnopqrstuvwxyzABCDEFGHIJKLMNOPQRSTUVWXYZ\relax}

% Notation of main.tex.
\newcommand{\ket}[1]{\left|#1\right\rangle}
\newcommand{\Pauli}{\mathcal P}
\newcommand{\Clifford}{\mathcal C}

% A font size in pt, for text matched to the submitted figures.
\newcommand{\figsize}[1]{\fontsize{#1}{1.2\dimexpr#1pt\relax}\selectfont}

% then the tikzpicture. Compile it from paper/figures/ (pdflatex <name>.tex);
% the PDF lands next to it and main.tex includes it as before.
%
% Only what the figures share lives here: fonts, packages, notation.
% Each figure keeps the colours of its submitted (Inkscape) version and
% defines them itself.
\usepackage[T1]{fontenc}
\usepackage{lmodern}
\usepackage{amsmath}
\usepackage{amssymb}
\usetikzlibrary{arrows.meta, calc, fit, backgrounds}

% Fonts. The submitted figures use one of two looks, so each figure picks
% its own: serif (Latin Modern, the default here) or, by calling
% \figsansfonts in its preamble, sans text with sans Latin letters and
% digits in math (as beamer does, in the scalable Latin Modern Sans),
% Greek left in cmmi.
\def\sfletter#1{\ifx#1\relax\else
  \DeclareMathSymbol{#1}{\mathalpha}{sfletters}{`#1}\expandafter\sfletter\fi}
\newcommand{\figsansfonts}{%
  \renewcommand{\familydefault}{\sfdefault}%
  \DeclareSymbolFont{sfletters}{OT1}{lmss}{m}{sl}%
  \SetSymbolFont{operators}{normal}{OT1}{lmss}{m}{n}%
  \sfletter abcdefghijklmnopqrstuvwxyzABCDEFGHIJKLMNOPQRSTUVWXYZ\relax}

% Notation of main.tex.
\newcommand{\ket}[1]{\left|#1\right\rangle}
\newcommand{\Pauli}{\mathcal P}
\newcommand{\Clifford}{\mathcal C}

% A font size in pt, for text matched to the submitted figures.
\newcommand{\figsize}[1]{\fontsize{#1}{1.2\dimexpr#1pt\relax}\selectfont}
